\paragraph{A comparison on the final spaces.}
Let \(f_i:E_i\to F\) be a finite nonempty family of integral maps of
nonzero Hermitian lattices over one number field \(K\). Use ordinary
slopes, normalized by \([K:\mathbb Q]\), and the map height
\[
 h(f)=\frac1{[K:\mathbb Q]}\sum_{\sigma:K\hookrightarrow\mathbb C}
                  \log\|f_\sigma\|,\qquad h(0)=-\infty.
\]
Write \(r_i=\operatorname{rank}E_i\), \(R=\sum_i r_i\), and
\(n=\operatorname{rank}F\). Suppose matrices
\(T_i:\mathbb C^{r_i}\to\mathbb C^n\) represent these same maps at
one complex embedding, in individual source bases and one common
target basis. All matrix norms and adjoints below use the auxiliary
Euclidean structures. Put
\[
 Q=\sum_iT_iT_i^*,\qquad M=\operatorname{tr}Q>0,\qquad
 \mathcal S=\bigoplus_iT_i^*T_i.
\]
For a self-adjoint operator \(A\), write \(A_+\) for its positive part.
Define
\[
 \gamma=\frac1M\operatorname{tr}\left(\frac MR I_R-\mathcal S\right)_+,
 \qquad
 \beta=\frac1M\operatorname{tr}\left(Q-\frac Mn I_n\right)_+,
\]
\[
 a^2=\frac R{M^2}\left\|\mathcal S-\frac MR I_R\right\|_{\rm HS}^2,
 \quad a\ge0.
\]
These quantities are unchanged by a common positive rescaling of the
matrices. For \(U,H\ge0\), set
\[
 L=U+H,\qquad
 \Phi(U,H)=
 \begin{cases}\sqrt{UH},&U\le H,\\(U+H)/2,&U>H,\end{cases}
 \qquad \mathcal A=\min\{a\Phi(U,H),\gamma L\}.
\]

\begin{theorem}[A slope bound from two trace errors]
\label{shared-trace-slope}
Let \(s,h_i\in\mathbb R\), and suppose
\[
 h(f_i)\le h_i,\qquad
 \frac1R\sum_i r_i\{\mu(E_i)-h_i\}\ge s,\qquad
 \mu_{\max}(E_i)-h_i\le s+U,\qquad
 \mu_{\min}(F)\ge s-H.
\]
Then
\begin{equation}\label{shared-trace-slope-bound}
 s-\mu(F)\le P_{\rm tr}:=\min\{H,\mathcal A+\beta L\}.
\end{equation}
In particular, \(\min\{H,(\gamma+\beta)(U+H)\}\) is a valid bound
using only the two trace errors.
\end{theorem}

\begin{proof}
Multiply the source norms by \(e^{h_i}\). Their slopes decrease by
\(h_i\), and all map heights become nonpositive. Such an integral
map preserves the Harder--Narasimhan filtration: a nonzero induced
map from slopes at least \(\xi\) to a quotient of slopes below
\(\xi\) would contradict the determinant inequality.

For \(\xi\in[s-H,s+U]\), let \(p(\xi)\) be the total source
filtration rank divided by \(R\), and let \(q(\xi)\) be the target
filtration rank divided by \(n\). Let \(P=\bigoplus_iP_i\) and \(P_W\)
be the corresponding Euclidean orthogonal projections. The image
inclusion and positivity give
\[
 \operatorname{tr}(P\mathcal S)
 =\operatorname{tr}\left(\sum_iT_iP_iT_i^*\right)
 \le\operatorname{tr}(P_WQ)\le M(q+\beta).
\]
The left side is at least \(M(p-\gamma)\). Centering \(\mathcal S\)
and using \(\|P-pI_R\|_{\rm HS}^2=Rp(1-p)\) also bounds it below by
\(M\{p-a\sqrt{p(1-p)}\}\).

The upper source slope bound and its mean give
\(\int p\ge H\) and \(z:=\int(1-p)\le U\) on this interval.
Consequently
\[
 \int\sqrt{p(1-p)}
 \le\sqrt{z(L-z)}\le\Phi(U,H).
\]
The lower target slope bound gives
\(\mu(F)\ge s-H+\int q\). Integrating either source lower bound
therefore proves \(s-\mu(F)\le\mathcal A+\beta L\).
The target minimum slope independently gives \(s-\mu(F)\le H\).
If \(L=0\), this last bound proves the result directly.
\end{proof}

For a trace-zero self-adjoint matrix, the positive and negative
traces both equal half its trace norm. Cauchy--Schwarz on the
eigenvalues therefore gives
\[
 \gamma\le a/2,\qquad
 \beta\le\frac{\sqrt\eta}{2},\qquad
 \eta=\frac n{M^2}\left\|Q-\frac Mn I_n\right\|_{\rm HS}^2.
\]
Thus the trace bound includes the centered squared-energy bound.
More generally, positive operators \(E_{\rm s},E_{\rm t}\) satisfying
\[
 \mathcal S\ge(M/R)I_R-E_{\rm s},\qquad
 Q\le(M/n)I_n+E_{\rm t}
\]
give \(\gamma\le\operatorname{tr}E_{\rm s}/M\) and
\(\beta\le\operatorname{tr}E_{\rm t}/M\).
Indeed, evaluate each inequality on the spectral projection defining
the corresponding positive part. No small operator norm is required.

\begin{corollary}[A direct upper period bound]
\label{shared-trace-bessel}
Under the same arithmetic hypotheses, suppose
\[
 Q\le(1+\lambda)(M/n)I_n+E,\qquad
 \lambda,\varepsilon\ge0,\quad E\ge0,\quad
 \operatorname{tr}E\le\varepsilon M.
\]
Then
\begin{equation}\label{shared-trace-bessel-bound}
 s-\mu(F)\le
 \min\left\{H,\frac{\mathcal A+\lambda H+\varepsilon L}
                         {1+\lambda}\right\}.
\end{equation}
For \(E=0\), take \(\varepsilon=0\). A supplied upper bound for
\(\lambda\) may replace it in this minimum.
\end{corollary}

\begin{proof}
The preceding rank argument now has upper bound
\(M\{(1+\lambda)q+\varepsilon\}\). Integration gives
\((1+\lambda)(s-\mu(F))\le\mathcal A+\lambda H+\varepsilon L\).
For the substitution assertion, write the fraction as
\(H+(\mathcal A+\varepsilon L-H)/(1+\lambda)\).
It increases with \(\lambda\) when its numerator after \(H\) is
nonpositive; otherwise the minimum equals \(H\).
\end{proof}

\begin{corollary}[An ambient trace comparison]\label{shared-trace-ambient}
Retain the arithmetic hypotheses of Theorem~\ref{shared-trace-slope}.
The matrices may instead have \(D\ge n\) rows, with their first \(n\)
rows representing the same maps in one common target basis. Define
\(M,\mathcal S,a,\gamma\) from these enlarged matrices, and define
\(\beta=M^{-1}\operatorname{tr}(Q-(M/D)I_D)_+\).
Suppose \(M>0\), and put \(\nu=n/D\), \(\delta=1-\nu\).
Then, with \(\mathcal A=\min\{a\Phi(U,H),\gamma(U+H)\}\),
\begin{equation}\label{shared-trace-ambient-bound}
 s-\mu(F)\le
 \min\left\{H,\frac{\delta U+\mathcal A+\beta(U+H)}{\nu}\right\}.
\end{equation}
\end{corollary}

\begin{proof}
Adjoin the \(D-n\) discarded coordinates to each target filtration
space. Its normalized dimension is \(\delta+\nu q\), so the previous
trace upper bound becomes \(M(\delta+\nu q+\beta)\).
Integrating either source lower bound gives
\(\nu\int q\ge H-\delta(U+H)-\mathcal A-\beta(U+H)\).
Use \(\mu(F)\ge s-H+\int q\) and \(\nu=1-\delta\).
The independent lower slope bound gives the minimum with \(H\).
\end{proof}

\paragraph{Restriction to the actual source and target.}
The next lemma starts with maps \(T_i:H_i\to H\) into an ambient
space of dimension \(D>0\). Define \(R,M,\mathcal S,\gamma\) as
above, and use \(D\) in the definition of \(\beta\).
Let \(J_i:H_i'\to H_i\) be isometries with
\(R'=\sum_i\dim H_i'>0\), and let \(\Pi\) be an orthogonal
target projection of rank \(n>0\). Remove zero-dimensional source
spaces when applying the arithmetic theorem. Put
\[
 T_{i,0}=\Pi T_iJ_i,\quad
 Q'=\sum_iT_iJ_iJ_i^*T_i^*,\quad
 Q_0=\Pi Q'\Pi|_{\operatorname{ran}\Pi},
\]
\[
 M_0=\operatorname{tr}Q_0,\quad
 \rho=\frac{\operatorname{tr}(Q-Q')}{M},\quad
 \delta=1-\frac nD .
\]

\begin{lemma}[Both restrictions measured by trace]
\label{shared-trace-restriction}
If \(k=1-\rho-\delta-\beta>0\), the final maps have positive mass
and their trace errors satisfy
\begin{equation}\label{shared-trace-restriction-bound}
 M_0\ge Mk,\qquad
 \beta_0\le\frac{\beta+\rho}{k},\qquad
 \gamma_0\le\frac{\gamma+\delta+\beta}{k},\qquad
 \rho\le\frac{R-R'}R+\gamma.
\end{equation}
Here \(\gamma_0,\beta_0\) use the actual dimensions \(R',n\)
and the actual mass \(M_0\).
Let \(a_0\ge0\) be the final centered source defect, defined by
\[
 a_0^2=\frac{R'}{M_0^2}
 \left\|\bigoplus_iT_{i,0}^*T_{i,0}-\frac{M_0}{R'}I_{R'}\right\|_{\rm HS}^2.
\]
If \(\Delta=\sum_i\|T_i^*T_i-I\|_{\rm HS}^2\), and \(\Delta_0\)
denotes the analogous final quantity, then
\[
 \Delta_0\le\Delta+2M(\delta+\beta),\qquad
 a_0^2=\frac{R'\Delta_0-(M_0-R')^2}{M_0^2}.
\]
\end{lemma}

\begin{proof}
If \(A\) is self-adjoint of trace zero and its compression \(B\)
acts on a space of dimension \(m>0\) and has trace \(b\), then
\(\operatorname{tr}(B-(b/m)I)_+\le\operatorname{tr}A_+\).
For \(b\ge0\), bound the left side by \(\operatorname{tr}B_+\);
for \(b\le0\), use the negative part instead.
Compression cannot increase either positive or negative trace,
by the variational formula
\(\operatorname{tr}A_+=\max_{0\le P\le I}\operatorname{tr}(PA)\).

Write \(L_{\rm s}=\rho M\). Source restriction changes the centered
target operator to
\[
 Q'-\frac{M-L_{\rm s}}D I
  =Q-\frac MD I-(Q-Q')+\frac{L_{\rm s}}D I.
\]
Its positive trace is at most \(M\beta+L_{\rm s}\).
Compression and recentering therefore give
\(\beta_0M_0\le M(\beta+\rho)\).
The target mass removed afterward is
\[
 L_{\rm t}=\operatorname{tr}((I-\Pi)Q')
            \le\operatorname{tr}((I-\Pi)Q)
            \le M(\delta+\beta).
\]
Hence \(M_0=M-L_{\rm s}-L_{\rm t}\ge Mk>0\).

Applying the same compression observation to the block diagonal
source operator shows that its unnormalized centered negative
trace after source restriction is at most \(M\gamma\).
Target restriction subtracts a positive block operator of trace
\(L_{\rm t}\). After recentering, its negative trace is therefore
at most \(M\gamma+L_{\rm t}\), proving the bound for \(\gamma_0\).
The trace of \(\mathcal S\) on the discarded source space is at most
\(M(R-R')/R+M\gamma\), proving the bound for \(\rho\).

Finally let \(G\) be the block source Gram operator after source
restriction, and let \(B\ge0\) be its discarded-target part.
Then \(0\le B\le G\), \(\operatorname{tr}B=L_{\rm t}\), and
\(\|G-I\|_{\rm HS}^2\le\Delta\). Expanding the square and using
\(\operatorname{tr}((G-B)B)\ge0\) gives
\(\|G-B-I\|_{\rm HS}^2\le\Delta+2L_{\rm t}\).
Centering this final operator gives the identity for \(a_0^2\).
\end{proof}

One may stop after source restriction and use the ambient comparison.
If \(\rho<1\), the matrices \(T_iJ_i\) have
\[
 M'=M(1-\rho)>0,\qquad
 \gamma'\le\frac{\gamma}{1-\rho},\qquad
 \beta'\le\frac{\beta+\rho}{1-\rho}.
\]
Suppose their first \(n\) rows represent the arithmetic maps in one
common target basis, and the slope and height hypotheses hold for
these restricted source lattices and that same target.
With \(\nu=n/D>0\), Corollary~\ref{shared-trace-ambient} gives
\begin{equation}\label{shared-trace-source-only}
 s-\mu(F)\le
 \min\left\{H,\frac{\delta U+
              \dfrac{\gamma+\beta+\rho}{1-\rho}(U+H)}{\nu}\right\}.
\end{equation}
The three estimates use only the source-compression steps of the
preceding proof. This route retains the target-rank term \(\delta U\)
without imposing \(k>0\) or estimating the final target mass.

\paragraph{The same spectral error in the source Gram matrix.}
\begin{lemma}[Transfer of the positive spectral trace]
\label{shared-trace-spectrum}
Let \(S:\bigoplus_iH_i\to H\) be synthesis,
\(S(v_i)=\sum_iT_iv_i\), and let
\(\mathcal K=S^*S=[T_i^*T_j]\). For every \(\alpha>0\),
\begin{equation}\label{shared-trace-spectrum-identity}
 \operatorname{tr}(Q-\alpha I_D)_+
       =\operatorname{tr}(\mathcal K-\alpha I_R)_+.
\end{equation}
Consequently a bound on the right by \(\varepsilon M\) supplies
\(Q\le\alpha I_D+E\) with \(E\ge0\) and
\(\operatorname{tr}E\le\varepsilon M\).
\end{lemma}

\begin{proof}
The positive eigenvalues of \(SS^*\) and \(S^*S\) agree, with
multiplicities: on a positive eigenspace, \(S\) and \(S^*\) give
inverse isomorphisms after division by its eigenvalue.
Any additional zero eigenvalues contribute zero to either positive
part. Take \(E=(Q-\alpha I_D)_+\).
\end{proof}

If an absolutely convergent block expansion gives
\(\mathcal K=I_R+N\), the required quantity is
\(\operatorname{tr}((1-\alpha)I_R+N)_+\).
Combine the signed terms before taking the positive part, retaining
every off-diagonal term unless its vanishing is proved.
For self-adjoint matrices \(A,B\), the variational formula gives
\[
 \left|\operatorname{tr}A_+-\operatorname{tr}B_+\right|
       \le R\|A-B\|_{\rm op}.
\]
Thus bounds for self-adjoint finite truncations pass to the limit
in operator norm; a quantitative truncation error retains the
factor \(R\). An expansion supplied only for the full period maps
can be used with Lemma~\ref{shared-trace-restriction}.
An expansion for already restricted maps must refer to those maps.
